\scanpage{221}
Soit $X$ \deleted{sur $S$} projectif sur $S$ de valuation discrète, définissant une famille de quadriques non singulières dégénérant en un cône quadratique
\[
X=\operatorname{Proj}\left(\frac{k[x,y,z,t]}{(x^2+y^2+z^2-\lambda t^2)}\right)\quad\text{sur }S=\operatorname{Spec}(\deleted{\unsure{K|k}}k[\lambda])
\]
\deleted{\unsure{$S$ anneau de valuation}}
% slanted
\margin{\unsure{pr l'utiliser immédiatement}, on a un \unsure{diviseur|fibré} sur $P^3_S$ \uline{\unsure{transversal}} aux fibres \ldots} $k$ un corps (si on ne le \unsure{dégrossit|dégoûte} pas, on prendrait $k=\mathbb Z$). $X$ est \deleted{\uline{séparable} sur $S$,} \inserted{\uline{$S$-\unsure{versel|normal|universel}} \uline{et} \uline{$S$-irréductible}}. \deleted{à fibres \unsure{abs.} irréductibles et} Pour calculer $\underline{\underline{\mathrm{Pic}}}_{X/S}$, il faut calculer $\mathrm{Pic}(X'/S')$ pour tout $S'$ \uline{local} sur $S$, i.e. $S'=\operatorname{Spec}(A)$, $A$ anneau \uline{local} noethérien muni d'un $\lambda\in\deleted{\unsure{r(A)}}A$. 1°) $\lambda\in\unsure{r(A)|\mathfrak r(A)|\mathfrak m(A)}$. Si $A$ est un corps, on trouve $\mathbb Z$, \inserted{générateur $\underline{\mathcal O}_{X'}(1)$,} on déduit que c'est encore vrai pour $A$ quelconque, grâce à $\tikz[baseline=(eq.base)]{\node[inner sep=2pt] (eq) {$H^1(X'_0,\underline{\mathcal O}_{X'_0})=0$};\draw (eq.center) ellipse (1.95cm and .34cm);}$.
% slanted
\margin{\uline{à mettre} \uline{\unsure{au lemme|ailleurs}}. Le schéma de Picard d'un cône étant $\mathbb Z$, son $H^1\deleted{\unsure{(\mathcal O_{X})}}$-additif est nul.} Si au contraire $\xi\notin r(A)$, alors la même méthode donne \deleted{$\unsure{H^2}$} $\mathbb Z^2$, de générateurs $\alpha,\beta$. D'où \deleted{l'image d'un} \inserted{le schéma de Picard,} qui a des parties fermées de type fini sur $S$ non propre sur $S$.
% archive sheet number: 220
